\documentclass[10pt,a4paper]{amsart}
\usepackage{amsmath,amssymb,amsthm,graphicx}
% Retained preamble for the Stacks Project chapter divisors.tex (Divisors).
% The chapter file pinned-original-source.tex carries only \input{preamble} on its first line, so no
% preamble text lives inside the chapter itself.  The Stacks Project uses one single shared preamble for
% every chapter; the snapshot stages it verbatim in sources/stacks-chow/preamble.tex.  The block below is
% copied mechanically, line for line, from that file's lines 14-18 (the Xy-pic setup) and lines 154-217
% (the native theorem declarations together with the mathematical macros).  Nothing was added, removed or
% redefined.  The wrapper class and the amsmath/amssymb/amsthm/graphicx base packages are supplied by the
% assembler, so the class/\externaldocument part of the shared preamble is not repeated here.
% For commutative diagrams we use Xy-pic
\usepackage[all]{xy}

% We use 2cell for 2-commutative diagrams.
\xyoption{2cell}
% Theorem environments.
%
\theoremstyle{plain}
\newtheorem{theorem}[subsection]{Theorem}
\newtheorem{proposition}[subsection]{Proposition}
\newtheorem{lemma}[subsection]{Lemma}

\theoremstyle{definition}
\newtheorem{definition}[subsection]{Definition}
\newtheorem{example}[subsection]{Example}
\newtheorem{exercise}[subsection]{Exercise}
\newtheorem{situation}[subsection]{Situation}

\theoremstyle{remark}
\newtheorem{remark}[subsection]{Remark}
\newtheorem{remarks}[subsection]{Remarks}

\numberwithin{equation}{subsection}

% Macros
%
\def\lim{\mathop{\mathrm{lim}}\nolimits}
\def\colim{\mathop{\mathrm{colim}}\nolimits}
\def\Spec{\mathop{\mathrm{Spec}}}
\def\Hom{\mathop{\mathrm{Hom}}\nolimits}
\def\Ext{\mathop{\mathrm{Ext}}\nolimits}
\def\SheafHom{\mathop{\mathcal{H}\!\mathit{om}}\nolimits}
\def\SheafExt{\mathop{\mathcal{E}\!\mathit{xt}}\nolimits}
\def\Sch{\mathit{Sch}}
\def\Mor{\mathop{\mathrm{Mor}}\nolimits}
\def\Ob{\mathop{\mathrm{Ob}}\nolimits}
\def\Sh{\mathop{\mathit{Sh}}\nolimits}
\def\NL{\mathop{N\!L}\nolimits}
\def\CH{\mathop{\mathrm{CH}}\nolimits}
\def\proetale{{pro\text{-}\acute{e}tale}}
\def\etale{{\acute{e}tale}}
\def\QCoh{\mathit{QCoh}}
\def\Ker{\mathop{\mathrm{Ker}}}
\def\Im{\mathop{\mathrm{Im}}}
\def\Coker{\mathop{\mathrm{Coker}}}
\def\Coim{\mathop{\mathrm{Coim}}}

% Boxtimes
%
\DeclareMathSymbol{\boxtimes}{\mathbin}{AMSa}{"02}

%
% Macros for moduli stacks/spaces
%
\def\QCohstack{\mathcal{QC}\!\mathit{oh}}
\def\Cohstack{\mathcal{C}\!\mathit{oh}}
\def\Spacesstack{\mathcal{S}\!\mathit{paces}}
\def\Quotfunctor{\mathrm{Quot}}
\def\Hilbfunctor{\mathrm{Hilb}}
\def\Curvesstack{\mathcal{C}\!\mathit{urves}}
\def\Polarizedstack{\mathcal{P}\!\mathit{olarized}}
\def\Complexesstack{\mathcal{C}\!\mathit{omplexes}}
% \Pic is the operator that assigns to X its picard group, usage \Pic(X)
% \Picardstack_{X/B} denotes the Picard stack of X over B
% \Picardfunctor_{X/B} denotes the Picard functor of X over B
\def\Pic{\mathop{\mathrm{Pic}}\nolimits}
\def\Picardstack{\mathcal{P}\!\mathit{ic}}
\def\Picardfunctor{\mathrm{Pic}}
\def\Deformationcategory{\mathcal{D}\!\mathit{ef}}

% --- appended 2026-10-12 (Chief): the project own definitions of the reference/slogan/history
% display environments, copied mechanically from the shared Stacks preamble lines 8-12
% (\usepackage{verbatim} plus the three \newenvironment lines, which the chapter preamble keeps
% verbatim from the project preamble). Nothing added, removed or redefined: these environments
% carry no mathematical content and the project itself typesets them as comments. Without these
% three lines a counted statement whose own text uses them fails with Environment slogan undefined.
\usepackage{verbatim}
\newenvironment{reference}{\comment}{\endcomment}
\newenvironment{slogan}{\comment}{\endcomment}
\newenvironment{history}{\comment}{\endcomment}
\usepackage{cleveref}
\crefname{originallemma}{original lemma}{original lemmas}
\Crefname{originallemma}{Original lemma}{Original lemmas}
\makeatletter
\newlabel{algebra-lemma-characterize-smooth-over-field}{{00TT}{source}{}{sourcepointer-1}{}}
\newlabel{algebra-lemma-characterize-smooth-over-field@cref}{{[originallemma][0][]00TT}{[1][0][]source}{}{}{}}
\makeatother
\crefname{originallemma}{original lemma}{original lemmas}
\Crefname{originallemma}{Original lemma}{Original lemmas}
\makeatletter
\newlabel{algebra-lemma-lci-local}{{00SF}{source}{}{sourcepointer-2}{}}
\newlabel{algebra-lemma-lci-local@cref}{{[originallemma][0][]00SF}{[1][0][]source}{}{}{}}
\makeatother
\crefname{originallemma}{original lemma}{original lemmas}
\Crefname{originallemma}{Original lemma}{Original lemmas}
\makeatletter
\newlabel{lemma-Noetherian-scheme-regular-ideal}{{063I}{source}{}{sourcepointer-3}{}}
\newlabel{lemma-Noetherian-scheme-regular-ideal@cref}{{[originallemma][0][]063I}{[1][0][]source}{}{}{}}
\makeatother
\crefname{originallemma}{original lemma}{original lemmas}
\Crefname{originallemma}{Original lemma}{Original lemmas}
\makeatletter
\newlabel{lemma-fibre-quasi-regular}{{063W}{source}{}{sourcepointer-4}{}}
\newlabel{lemma-fibre-quasi-regular@cref}{{[originallemma][0][]063W}{[1][0][]source}{}{}{}}
\makeatother
\crefname{originallemma}{original lemma}{original lemmas}
\Crefname{originallemma}{Original lemma}{Original lemmas}
\makeatletter
\newlabel{morphisms-lemma-finite-type-noetherian}{{01T6}{source}{}{sourcepointer-5}{}}
\newlabel{morphisms-lemma-finite-type-noetherian@cref}{{[originallemma][0][]01T6}{[1][0][]source}{}{}{}}
\makeatother
\crefname{originallemma}{original lemma}{original lemmas}
\Crefname{originallemma}{Original lemma}{Original lemmas}
\makeatletter
\newlabel{morphisms-lemma-local-complete-intersection}{{01UG}{source}{}{sourcepointer-6}{}}
\newlabel{morphisms-lemma-local-complete-intersection@cref}{{[originallemma][0][]01UG}{[1][0][]source}{}{}{}}
\makeatother
\crefname{originallemma}{original lemma}{original lemmas}
\Crefname{originallemma}{Original lemma}{Original lemmas}
\makeatletter
\newlabel{morphisms-lemma-smooth-flat}{{01VF}{source}{}{sourcepointer-7}{}}
\newlabel{morphisms-lemma-smooth-flat@cref}{{[originallemma][0][]01VF}{[1][0][]source}{}{}{}}
\makeatother
\crefname{originallemma}{original lemma}{original lemmas}
\Crefname{originallemma}{Original lemma}{Original lemmas}
\makeatletter
\newlabel{morphisms-lemma-smooth-locally-finite-presentation}{{01VE}{source}{}{sourcepointer-8}{}}
\newlabel{morphisms-lemma-smooth-locally-finite-presentation@cref}{{[originallemma][0][]01VE}{[1][0][]source}{}{}{}}
\makeatother
\crefname{originallemma}{original lemma}{original lemmas}
\Crefname{originallemma}{Original lemma}{Original lemmas}
\makeatletter
\newlabel{morphisms-lemma-syntomic-flat}{{01UL}{source}{}{sourcepointer-9}{}}
\newlabel{morphisms-lemma-syntomic-flat@cref}{{[originallemma][0][]01UL}{[1][0][]source}{}{}{}}
\makeatother
\crefname{originallemma}{original lemma}{original lemmas}
\Crefname{originallemma}{Original lemma}{Original lemmas}
\makeatletter
\newlabel{morphisms-lemma-syntomic-locally-finite-presentation}{{01UK}{source}{}{sourcepointer-10}{}}
\newlabel{morphisms-lemma-syntomic-locally-finite-presentation@cref}{{[originallemma][0][]01UK}{[1][0][]source}{}{}{}}
\makeatother
\title{An immersion of a syntomic scheme into a smooth scheme over the base is a regular immersion}
\author{The Stacks Project Authors}
\date{stacks-project:divisors.tex:local-snapshot-20261012; GNU FDL 1.2 or later}
\begin{document}
\maketitle
\noindent\textit{Editorial note.} All delivered mathematical text is the author's own, copied from the pinned Stacks Project chapter file pinned-original-source.tex; no statement or proof step was rewritten, reordered or removed.\par\medskip
\noindent\textit{Editorial note.} Following the multi-claim convention for this source, source\_url is the source-level URL https://stacks.math.columbia.edu/ and the source label lemma-immersion-lci-into-smooth-regular-immersion is carried in claim\_key and here rather than in a per-claim tag URL.\par\medskip
\noindent\textit{Editorial note.} The preamble is the chapter's shared Stacks Project preamble, retained as sources/stacks-divisors/retained-preamble.tex: a line-for-line excerpt of the shared Stacks preamble staged with this snapshot. No macro was invented or redefined, and this card adds no preamble text of its own.\par\medskip
\noindent\textit{Editorial note.} Statement span 5672-5684 and proof span 5686-5717 come from the halfA pool card operations/pool-cards/stacks-divisors-halfA.json and were re-read in the pinned chapter file before this card was written.\par\medskip
\noindent\textit{Editorial note.} The statement contains an xymatrix commutative triangle; it is copied verbatim and the retained chapter preamble loads Xy-pic.\par\medskip
\noindent\textit{Editorial note.} 10 cross-references occur inside the counted spans; they are listed in external\_references as navigation pointers only.\par\medskip
\noindent\textit{Allowed standard prerequisites (editorial).}\begin{itemize}
\item A syntomic morphism is flat and locally of finite presentation, and a smooth morphism is flat and locally of finite presentation, cited from the morphisms chapter. Cited at proof lines 5686-5692. Supplies the hypotheses of the fibrewise criterion for a relative quasi-regular immersion.
\item The fibrewise criterion producing a relative quasi-regular immersion from a closed subscheme whose fibres are quasi-regular, cited from the same chapter. Cited at proof lines 5692-5694. Reduces the claim to the statement that each fibre is a quasi-regular closed subscheme.
\item Over a locally Noetherian scheme regularity of an immersion can be checked at points, cited from the same chapter, and a finite type scheme over a field is Noetherian. Cited at proof lines 5694-5699. Allows the fibre statement to be checked at individual points.
\item The local ring of a point of a fibre of a syntomic morphism is a local complete intersection over the residue field of the base point, the local ring of a point of a smooth scheme over a field is regular, and a quotient of a regular local ring by an ideal defining a local complete intersection is cut out by a regular sequence. Cited at proof lines 5701-5711. Provides the pointwise algebra input that closes the argument.
\end{itemize}\medskip
\ifdefined\hypertarget\hypertarget{sourcepointer-1}{}\fi
\noindent\textit{Source reference.} Original lemma 00TT: focused-seeds/AGP/sources/stacks-snapshot-20261011/tags.txt, source lines 1087--1087. Same-chapter reference used in the retained author text at source line(s) 5629, 5711. It is supplied only as a navigation pointer so that the retained text compiles; the counted item delivers the author's own statement and the author's complete proof, and no part of the referenced item is copied in or substituted for the counted argument.\par\medskip
\ifdefined\hypertarget\hypertarget{sourcepointer-2}{}\fi
\noindent\textit{Source reference.} Original lemma 00SF: focused-seeds/AGP/sources/stacks-snapshot-20261011/tags.txt, source lines 1037--1037. Same-chapter reference used in the retained author text at source line(s) 5713. It is supplied only as a navigation pointer so that the retained text compiles; the counted item delivers the author's own statement and the author's complete proof, and no part of the referenced item is copied in or substituted for the counted argument.\par\medskip
\ifdefined\hypertarget\hypertarget{sourcepointer-3}{}\fi
\noindent\textit{Source reference.} Original lemma 063I: focused-seeds/AGP/sources/stacks-snapshot-20261011/tags.txt, source lines 7832--7832. Same-chapter reference used in the retained author text at source line(s) 4939, 4970, 5169, 5234, 5705. It is supplied only as a navigation pointer so that the retained text compiles; the counted item delivers the author's own statement and the author's complete proof, and no part of the referenced item is copied in or substituted for the counted argument.\par\medskip
\ifdefined\hypertarget\hypertarget{sourcepointer-4}{}\fi
\noindent\textit{Source reference.} Original lemma 063W: focused-seeds/AGP/sources/stacks-snapshot-20261011/tags.txt, source lines 7845--7845. Same-chapter reference used in the retained author text at source line(s) 5619, 5699. It is supplied only as a navigation pointer so that the retained text compiles; the counted item delivers the author's own statement and the author's complete proof, and no part of the referenced item is copied in or substituted for the counted argument.\par\medskip
\ifdefined\hypertarget\hypertarget{sourcepointer-5}{}\fi
\noindent\textit{Source reference.} Original lemma 01T6: focused-seeds/AGP/sources/stacks-snapshot-20261011/tags.txt, source lines 2363--2363. Same-chapter reference used in the retained author text at source line(s) 1147, 5702. It is supplied only as a navigation pointer so that the retained text compiles; the counted item delivers the author's own statement and the author's complete proof, and no part of the referenced item is copied in or substituted for the counted argument.\par\medskip
\ifdefined\hypertarget\hypertarget{sourcepointer-6}{}\fi
\noindent\textit{Source reference.} Original lemma 01UG: focused-seeds/AGP/sources/stacks-snapshot-20261011/tags.txt, source lines 2409--2409. Same-chapter reference used in the retained author text at source line(s) 5707. It is supplied only as a navigation pointer so that the retained text compiles; the counted item delivers the author's own statement and the author's complete proof, and no part of the referenced item is copied in or substituted for the counted argument.\par\medskip
\ifdefined\hypertarget\hypertarget{sourcepointer-7}{}\fi
\noindent\textit{Source reference.} Original lemma 01VF: focused-seeds/AGP/sources/stacks-snapshot-20261011/tags.txt, source lines 2444--2444. Same-chapter reference used in the retained author text at source line(s) 5616, 5696. It is supplied only as a navigation pointer so that the retained text compiles; the counted item delivers the author's own statement and the author's complete proof, and no part of the referenced item is copied in or substituted for the counted argument.\par\medskip
\ifdefined\hypertarget\hypertarget{sourcepointer-8}{}\fi
\noindent\textit{Source reference.} Original lemma 01VE: focused-seeds/AGP/sources/stacks-snapshot-20261011/tags.txt, source lines 2443--2443. Same-chapter reference used in the retained author text at source line(s) 5617, 5697. It is supplied only as a navigation pointer so that the retained text compiles; the counted item delivers the author's own statement and the author's complete proof, and no part of the referenced item is copied in or substituted for the counted argument.\par\medskip
\ifdefined\hypertarget\hypertarget{sourcepointer-9}{}\fi
\noindent\textit{Source reference.} Original lemma 01UL: focused-seeds/AGP/sources/stacks-snapshot-20261011/tags.txt, source lines 2414--2414. Same-chapter reference used in the retained author text at source line(s) 5693. It is supplied only as a navigation pointer so that the retained text compiles; the counted item delivers the author's own statement and the author's complete proof, and no part of the referenced item is copied in or substituted for the counted argument.\par\medskip
\ifdefined\hypertarget\hypertarget{sourcepointer-10}{}\fi
\noindent\textit{Source reference.} Original lemma 01UK: focused-seeds/AGP/sources/stacks-snapshot-20261011/tags.txt, source lines 2413--2413. Same-chapter reference used in the retained author text at source line(s) 5692. It is supplied only as a navigation pointer so that the retained text compiles; the counted item delivers the author's own statement and the author's complete proof, and no part of the referenced item is copied in or substituted for the counted argument.\par\medskip
\begin{lemma}
\label{lemma-immersion-lci-into-smooth-regular-immersion}
Let
$$
\xymatrix{
Y \ar[rd] \ar[rr]_i & & X \ar[ld] \\
& S
}
$$
be a commutative diagram of morphisms of schemes.
Assume that $Y \to S$ is syntomic, $X \to S$ smooth, and
$i$ an immersion. Then $i$ is a regular immersion.
\end{lemma}

\begin{proof}
After replacing $X$ by an open neighbourhood of $i(Y)$
we may assume that $i$ is a closed immersion.
Let $T = i(Y)$ be the corresponding closed subscheme of $X$. Since
$T \cong Y$ the morphism $T \to S$ is flat and of finite presentation
(Morphisms, Lemmas
\ref{morphisms-lemma-syntomic-locally-finite-presentation} and
\ref{morphisms-lemma-syntomic-flat}).
Also a smooth morphism is flat and locally of finite presentation
(Morphisms, Lemmas
\ref{morphisms-lemma-smooth-flat} and
\ref{morphisms-lemma-smooth-locally-finite-presentation}).
Thus, according to
Lemma \ref{lemma-fibre-quasi-regular},
it suffices to show that $T_s \subset X_s$ is a quasi-regular closed
subscheme. As $X_s$ is locally of finite type over a field, it is Noetherian
(Morphisms, Lemma \ref{morphisms-lemma-finite-type-noetherian}).
Thus we can check that $T_s \subset X_s$ is a quasi-regular immersion
at points, see
Lemma \ref{lemma-Noetherian-scheme-regular-ideal}.
Take $t \in T_s$. By
Morphisms, Lemma \ref{morphisms-lemma-local-complete-intersection}
the local ring $\mathcal{O}_{T_s, t}$ is a local complete intersection
over $\kappa(s)$.
The local ring $\mathcal{O}_{X_s, t}$ is regular, see
Algebra, Lemma \ref{algebra-lemma-characterize-smooth-over-field}.
By
Algebra, Lemma \ref{algebra-lemma-lci-local}
we see that the kernel of the surjection
$\mathcal{O}_{X_s, t} \to \mathcal{O}_{T_s, t}$ is generated by a regular
sequence, which is what we had to show.
\end{proof}

\end{document}
